\documentclass{article}
\usepackage{/globals/tables}
\usepackage{/globals/code}
\usepackage{/globals/anki}
\usepackage{/globals/de}
\ankiprefix{std}
\ankitopic{STD}
\begin{document}
\section{STD} 
Java hat eine ausführliche Standardlibrary; Klassen und Methoden für nützliche, häufig genutzte Aufgaben, die bereits eingebaut sind.

\begin{ztable}{llX}{Klasse & Methode & Rückgabe}
\ilsourcecode|Math| & \ilsourcecode|abs(x)| & Den absoluten Wert von x, $\vert x \vert$ \\
& \ilsourcecode|min(x)| & Den kleineren von zwei Werten \\
& \ilsourcecode|max(x)| & Den größeren von zwei Werten \\
& \ilsourcecode|floor(x)| & Die abgerundete Zahl \\
& \ilsourcecode|ceil(x)| & Die aufgerundete Zahl \\
& \ilsourcecode|round(x)| & Die kaufmännisch gerundete Zahl \\
& \ilsourcecode|sqrt(x)| & Die Quadratwurzel \\
& \ilsourcecode|pow(b,a)| & Die Potenz $b^a$ \\
& \ilsourcecode|log(x)| & Den natürlichen Logarithmus $\log_e{x}$ \\
& \ilsourcecode|log10(x)| & Den dekadischen Logarithmus $\log_{10}{x}$ \\
& \ilsourcecode|random()| & Eine pseudozufallige Dezimalzahl, $[0;1[$ \\
& \ilsourcecode|sin(x)| & Den Sinus, $x$ in Rad \\
& \ilsourcecode|cos(x)| & Den Cosinus, $x$ in Rad \\
& \ilsourcecode|tan(x)| & Den Tangents, $x$ in Rad \\
& \ilsourcecode|asin(x)| & Den Arkussinus, $x$ in Rad \\
& \ilsourcecode|acos(x)| & Den Arkuscosinus, $x$ in Rad \\
& \ilsourcecode|atan(x)| & Den Arkustangents, $x$ in Rad \\
& \ilsourcecode|toRadians(x)| & Rechnet den Deg-Winkel $x$ in Rad um \\
& \ilsourcecode|toDegrees(x)| & Rechnet den Rad-Winkel $x$ in Deg um \\
\hline
\ilsourcecode|Integer| & \ilsourcecode|parseInt(s)| & Verwandelt einen String in ein int \\
\ilsourcecode|Double| & \ilsourcecode|parseDouble(s)| & Verwandelt einen String in ein double \\
\ilsourcecode|Boolean| & \ilsourcecode|parseBool(s)| & Verwandelt einen String in ein bool \\
\hline
\ilsourcecode|Arrays| & \ilsourcecode|toString(arr)| & Verwandelt ein Array in eine menschlich lesbare Darstellung \\
 & \ilsourcecode|equals(arr1,arr2)| & Überprüft ob zwei Arrays identisch sind \\
 & \ilsourcecode|copyOf(arr,n)| & Kopiert ein Array, $[n;]$ \\
 & \ilsourcecode|sort(arr)| & Sortiert ein Array \\
\end{ztable}
\ankinote{arc}{Arkus-}{a}
\ankinote{arreq}{Arrays sind gleich}{Arrays.equals(arr1,arr2)}
\ankinote{arrcopy}{Array kopieren}{Arrays.copyOf(arr,n)}
\ankinote{arrsort}{Array sortieren}{Arrays.sort(arr)}

\end{document}